\documentclass[11pt,a4paper]{article}
\usepackage{estilo-nexasoft}

\begin{document}
\hypersetup{pageanchor=false}

% ============================================================
% PORTADA
% ============================================================

\begin{titlepage}
\thispagestyle{empty}

\begin{tikzpicture}[remember picture,overlay]
\fill[Navy]
 (current page.south west)
 rectangle (current page.north east);

\fill[Teal]
 ([xshift=-17mm]current page.north east)
 rectangle
 ([xshift=-11mm]current page.south east);

\draw[white,opacity=.13,line width=1.1pt]
 ([xshift=-77mm,yshift=-50mm]current page.north east)
 circle (64mm);

\fill[Teal]
 ([xshift=23mm,yshift=-61mm]current page.north west)
 rectangle ++(28mm,2mm);
\end{tikzpicture}

\color{white}
\vspace*{15mm}

{\Large\bfseries NEXASOFT\par}
\vspace{2mm}
{\small SOFTWARE INSTITUCIONAL\par}

\vspace{49mm}

{\fontsize{37}{40}\selectfont\bfseries
Manual de\par usuario\par}

\vspace{14mm}

{\Large
Sistema de control de ingreso a exámenes\par}

\vspace{5mm}

{\small
GUÍA ILUSTRADA DE OPERACIÓN POR PERFILES\par}

\vfill

{\normalsize Universidad Mayor de San Simón\par}
{\small Taller de Ingeniería de Software\par}

\vspace{15mm}

{\small
Gestión 2/2026\qquad
Sprint 5\qquad
Edición de estabilización\par}

\end{titlepage}

\hypersetup{pageanchor=true}
\pagenumbering{roman}

% ============================================================
% CONTROL DOCUMENTAL
% ============================================================

\section*{Ficha de publicación}

\begin{tabularx}{\linewidth}{
 >{\bfseries}p{34mm}X
}
\toprule
Producto &
Punku — Sistema de control de ingreso a exámenes masivos\\
Elaborado por & NexaSoft S.R.L.\\
Destino & Personal administrativo, docente y auxiliar\\
Versión &
Manual de usuario — edición de estabilización, gestión 2/2026\\
Evidencia &
32 capturas locales: 16 pantallas en escritorio y 16 en móvil\\
\bottomrule
\end{tabularx}

\begin{notebox}
\textbf{Cómo leer este manual.}
Cada módulo reúne su propósito, la secuencia de uso,
la comprobación de resultados y dos vistas verificadas.
La captura de escritorio se reproduce completa; en móvil
se presenta el primer tramo visible del navegador para
mantener su legibilidad. Los 32 archivos PNG originales
permanecen íntegros en
\texttt{docs/manuales/imagenes/}.
\end{notebox}

\tableofcontents
\clearpage
\pagenumbering{arabic}

% ============================================================
\section{Antes de comenzar}
% ============================================================

\textbf{Acceso institucional:}
\url{https://nexasoft.tis.cs.umss.edu.bo/login}.

El menú depende de los permisos y asignaciones efectivos;
dos cuentas del mismo rol pueden ver opciones distintas.
Para comenzar, escriba el identificador de usuario y su
contraseña y confirme el acceso. Tras finalizar, utilice
\textit{Cerrar sesión}.

\begin{notebox}
\textbf{Protección de datos.}
No comparta contraseñas, códigos QR de acceso ni capturas
con datos personales. Las ilustraciones proceden de una
instalación local y no constituyen instrucciones para
modificar registros reales sin autorización.
\end{notebox}

% ============================================================
\section{Mapa de tareas por perfil}
% ============================================================

\begin{tabularx}{\linewidth}{>{\bfseries}p{30mm}X}
\toprule
Administrador &
Preparar período, oferta, aulas, docentes y padrón;
consultar reportes; administrar cuentas, bitácora y respaldos.\\
Docente &
Localizar exámenes, comprobar grupos, consultar reportes
propios, seguimiento y punto de control.\\
Auxiliar &
Consultar el punto de control y la lista de aula asignada.\\
Estudiante &
Dispone de un acceso propio; sus pantallas no forman parte
de las 32 capturas recopiladas.\\
\bottomrule
\end{tabularx}

% ============================================================
\section{Administrador: preparación y supervisión}
% ============================================================

\modulo{Períodos académicos}
{Administrador}{/periodo}{periodo}

Aquí se establece el contexto temporal de la actividad
académica. La vista muestra períodos detectados, pendientes
y la situación de la importación por facultad; no confunda
un período visible con uno necesariamente vigente.

\pasos
\begin{enumerate}
\item Acceda a \textbf{Período} desde el menú de preparación.
\item Localice el período por su etiqueta y fechas.
\item Compruebe el indicador de estado y, cuando corresponda,
los pendientes de docentes y grupos.
\item Abra el detalle o la opción de importación únicamente
si su cuenta está autorizada.
\end{enumerate}

\resultado{
Se identifican el período de trabajo, su estado
y las tareas de preparación pendientes.
}

\Pantallas
{administrador/escritorio/periodo.png}
{administrador/movil/periodo.png}
{853}
{Período académico: información global en escritorio
y navegación vertical en móvil.}

% ------------------------------------------------------------

\modulo{Oferta académica}
{Administrador}{/oferta}{oferta}

La sección organiza las asignaturas y sus grupos.
Los filtros de período, facultad, carrera, nivel y búsqueda
permiten acotar una lista extensa sin cambiar su contenido.

\pasos
\begin{enumerate}
\item Abra \textbf{Oferta}.
\item Seleccione los criterios disponibles de período
y unidad académica.
\item Escriba una asignatura o código en el buscador
si necesita localizar un registro.
\item Despliegue la asignatura para comprobar sus grupos.
\item Verifique que los filtros activos correspondan
a la consulta antes de continuar.
\end{enumerate}

\resultado{
La lista refleja el subconjunto de asignaturas y grupos
coincidente con los filtros seleccionados.
}

\Pantallas
{administrador/escritorio/oferta.png}
{administrador/movil/oferta.png}
{826}
{Oferta académica: filtros y listado de asignaturas
adaptados al dispositivo.}

% ------------------------------------------------------------

\modulo{Aulas y mapa}
{Administrador}{/aulas}{aulas}

El módulo permite identificar ambientes por campus,
edificio y aula. Su mapa es una referencia visual que
complementa el listado, no sustituye la verificación
del ambiente asignado al examen.

\pasos
\begin{enumerate}
\item Abra \textbf{Aulas y mapa}.
\item Seleccione el campus o la unidad académica cuando
el filtro esté disponible.
\item Busque el edificio y revise los ambientes asociados.
\item Contraste el aula seleccionada con su ubicación
en el mapa.
\end{enumerate}

\resultado{
Se puede relacionar el nombre del aula con el edificio
y la ubicación mostrada.
}

\Pantallas
{administrador/escritorio/aulas.png}
{administrador/movil/aulas.png}
{830}
{Aulas: mapa del campus y panel de ambientes;
en móvil los controles se apilan.}

% ------------------------------------------------------------

\modulo{Docentes}
{Administrador}{/docentes}{docentes}

La consulta docente permite revisar el directorio académico
y detectar registros que requieren atención. Los indicadores
de resumen y los filtros no deben confundirse con la
autorización para editar datos personales.

\pasos
\begin{enumerate}
\item Entre a \textbf{Docentes}.
\item Seleccione facultad o criterios de búsqueda
para restringir los resultados.
\item Revise el estado del registro localizado.
\item Consulte la información académica o de cuenta
únicamente si su permiso lo permite.
\end{enumerate}

\resultado{
El registro se encuentra sin confundir personas con nombres
similares y se visualiza su estado.
}

\Pantallas
{administrador/escritorio/docentes.png}
{administrador/movil/docentes.png}
{1197}
{Directorio docente y filtros; las tarjetas de resumen
permiten identificar registros pendientes.}

% ------------------------------------------------------------

\modulo{Padrón estudiantil}
{Administrador}{/estudiantes}{padron}

La vista del padrón contiene indicadores de población
y una lista de estudiantes. Una búsqueda debe verificarse
tanto por el criterio introducido como por los resultados,
especialmente si se gestionan registros masivos.

\pasos
\begin{enumerate}
\item Abra \textbf{Padrón}.
\item Seleccione facultad, carrera o estado según
los filtros disponibles.
\item Busque el estudiante mediante el identificador permitido.
\item Compruebe que la lista y los contadores correspondan
a la consulta.
\item Para cualquier carga masiva, use únicamente
el procedimiento autorizado y revise las validaciones
de archivo.
\end{enumerate}

\resultado{
El resultado corresponde a los filtros y no se asume
que una búsqueda sin coincidencias implique ausencia
definitiva del estudiante.
}

\nota{
Los registros estudiantiles son datos personales.
Evite divulgar listados obtenidos de esta sección.
}

\Pantallas
{administrador/escritorio/padron.png}
{administrador/movil/padron.png}
{2570}
{Padrón estudiantil: indicadores, búsqueda
y estado de los registros.}

% ------------------------------------------------------------

\modulo{Reportes institucionales}
{Administrador}{/reportes}{reportes}

Este módulo permite comparar exámenes por período,
facultad, texto de búsqueda y rango de fechas.
Las cifras agregadas y los archivos exportados
deben leerse con los mismos filtros activos.

\pasos
\begin{enumerate}
\item Abra \textbf{Reportes}.
\item Confirme el período y aplique los filtros de facultad,
búsqueda o fechas.
\item Revise los indicadores de ingresos, estudiantes
registrados e intentos rechazados.
\item Seleccione un examen para consultar su detalle
o bitácora.
\item Elija el reporte requerido y el formato habilitado
(PDF o XLSX).
\item Compruebe el período, filtros y contenido
del archivo descargado.
\end{enumerate}

\resultado{
Los indicadores y la exportación se interpretan con
el mismo alcance de consulta.
}

\Pantallas
{administrador/escritorio/reportes.png}
{administrador/movil/reportes.png}
{1369}
{Reportes institucionales: filtros, indicadores,
listado de exámenes y exportaciones.}

% ------------------------------------------------------------

\modulo{Usuarios, roles y permisos}
{Administrador}{/admin}{usuarios}

La sección muestra cuentas y una matriz de permisos
por rol. Un permiso modifica el alcance de acceso:
revisar la matriz es distinto de realizar una
modificación de seguridad.

\pasos
\begin{enumerate}
\item Abra \textbf{Usuarios y roles}.
\item Utilice búsqueda o filtro de rol para localizar
una cuenta.
\item Compruebe su estado y el rol asignado.
\item Examine la matriz de permisos y el alcance
de cada columna.
\item Solo si está autorizado, revise los cambios
antes de guardarlos y confirme su efecto.
\end{enumerate}

\resultado{
Se identifica quién puede acceder a una función
y con qué rol, sin asumir permisos por el título
del puesto.
}

\nota{
Los cambios de permisos son sensibles.
Deben coordinarse con la política de acceso institucional.
}

\Pantallas
{administrador/escritorio/usuarios-roles.png}
{administrador/movil/usuarios-roles.png}
{2678}
{Administración: listado de cuentas
y matriz de permisos por rol.}

% ------------------------------------------------------------

\modulo{Bitácora de operaciones}
{Administrador}{/bitacora}{bitacora}

La bitácora reúne eventos de auditoría. Su utilidad
depende de consultar el actor, la operación y el
momento correcto, preservando la integridad de los registros.

\pasos
\begin{enumerate}
\item Abra \textbf{Bitácora}.
\item Acote la búsqueda por usuario, operación y fechas
cuando esos filtros estén disponibles.
\item Examine el evento relacionado con la consulta
o incidencia.
\item Documente el hallazgo sin alterar ni borrar
el registro original.
\end{enumerate}

\resultado{
Los eventos consultados quedan identificados
por sus datos de auditoría.
}

\Pantallas
{administrador/escritorio/bitacora.png}
{administrador/movil/bitacora.png}
{1648}
{Bitácora: filtros de auditoría
y registros de operaciones.}

% ------------------------------------------------------------

\modulo{Respaldos}
{Administrador}{/respaldo}{respaldos}

La pantalla permite consultar el estado de los respaldos
y distinguir los programados, manuales o fallidos cuando
existan. Generar un respaldo no equivale a comprobar que
su restauración sea viable.

\pasos
\begin{enumerate}
\item Abra \textbf{Respaldo}.
\item Revise el estado y el historial disponible.
\item Identifique la última ejecución y cualquier
fallo registrado.
\item Si necesita un respaldo manual, solicite o ejecute
la operación únicamente con el permiso requerido.
\item Para una restauración, siga un procedimiento
técnico previamente verificado.
\end{enumerate}

\resultado{
Se reconoce el estado de los respaldos y se identifican
situaciones que requieren intervención.
}

\nota{
Nunca restaure datos en el entorno institucional
para probar una captura o un procedimiento del manual.
}

\Pantallas
{administrador/escritorio/respaldos.png}
{administrador/movil/respaldos.png}
{0}
{Respaldo: estado del módulo; las vistas de ejemplo
pueden mostrar ausencia de registros.}

% ============================================================
\section{Docente: preparación y control de exámenes}
% ============================================================

\modulo{Exámenes}{Docente}{/examenes}{examenes}

El panel organiza las evaluaciones por período y muestra
estados de preparación como grupos, aulas, personal
responsable, habilitación y códigos QR.
No todos los exámenes están listos para iniciar.

\pasos
\begin{enumerate}
\item Ingrese con la cuenta docente y abra \textbf{Exámenes}.
\item Localice la evaluación por fecha, asignatura y estado.
\item Revise los indicadores de grupos, aulas, quién atiende,
habilitación y QR.
\item Abra la acción disponible en la tarjeta,
según el estado del examen.
\item Antes de ingresar al control, confirme
que la preparación está completa.
\end{enumerate}

\resultado{
Cada tarjeta permite identificar el estado del examen
y la siguiente operación disponible.
}

\Pantallas
{docente/escritorio/examenes.png}
{docente/movil/examenes.png}
{1498}
{Panel docente: exámenes programados
y sus condiciones de preparación.}

% ------------------------------------------------------------

% ============================================================
% REGISTRO DE EXÁMENES
% Procedimiento funcional basado en RegistrarExamen.jsx
% Evidencias: cuatro etapas, escritorio y móvil
% ============================================================

\subsection{Registrar un examen: procedimiento completo}
\label{sec:registro-examen}

El registro de un examen se realiza desde
\textbf{Exámenes} mediante la opción
\textbf{Registrar examen}. La aplicación organiza el
procedimiento en cinco etapas: \textbf{Examen},
\textbf{Grupos}, \textbf{Aulas},
\textbf{Quién atiende} y \textbf{Listo}.

Las primeras cuatro etapas permiten preparar los datos.
La última muestra el resultado después de guardar
correctamente el examen.

\medskip
\begin{center}
\small\bfseries\color{Teal}
01 Examen
\quad--\quad
02 Grupos
\quad--\quad
03 Aulas
\quad--\quad
04 Personal
\quad--\quad
05 Listo
\end{center}
\medskip

\noindent\textbf{Antes de empezar.}
La cuenta debe tener permiso para gestionar exámenes,
disponer de una asignatura habilitada y trabajar
dentro del período académico correspondiente.

\noindent\textbf{Importante.}
La disponibilidad de asignaturas, grupos, ambientes
y personal depende de la información académica
registrada en el sistema.

% ============================================================
% MACRO EDITORIAL PARA LOS PARES DE CAPTURAS
% ============================================================

\newcommand{\CapturasRegistro}[3]{%
    \par\medskip
    \noindent
    \begin{minipage}{\linewidth}
        \centering
        \includegraphics[
            width=0.98\linewidth,
            height=0.39\textheight,
            keepaspectratio
        ]{imagenes/docente/flujo-registro/escritorio/paso-#1.png}
        \captionof{figure}{#2: vista de escritorio.}
        \label{fig:registro-escritorio-#1}
    \end{minipage}

    \par\medskip

    \noindent
    \begin{minipage}[t]{0.42\linewidth}
        \vspace{0pt}
        \centering
        \includegraphics[
            width=\linewidth,
            height=0.40\textheight,
            keepaspectratio
        ]{imagenes/docente/flujo-registro/movil/paso-#1.png}
        \captionof{figure}{#2: vista móvil.}
        \label{fig:registro-movil-#1}
    \end{minipage}
    \hfill
    \begin{minipage}[t]{0.53\linewidth}
        \vspace{3pt}
        {\bfseries\color{Teal}Uso desde el teléfono}\par
        \small #3
    \end{minipage}

    \par\medskip
}

% ============================================================
% ETAPA 1
% ============================================================

\subsubsection{Etapa 1. Datos generales del examen}

La primera pantalla determina qué evaluación se
realizará, a qué asignatura pertenece y cuándo tendrá
lugar. Estos datos condicionan las opciones de las
etapas posteriores.

\noindent\textbf{Procedimiento}

\begin{enumerate}
    \item Seleccione la \textbf{asignatura}
    que corresponde al examen.

    \item Elija el \textbf{tipo de evaluación}:
    primer parcial, segundo parcial, final,
    segunda instancia o mesa, según corresponda.

    \item Indique la \textbf{fecha}.
    Compruebe que pertenece al período y al
    calendario permitido para ese tipo de examen.

    \item Establezca la \textbf{hora de inicio}
    y la \textbf{duración} en minutos.

    \item Revise las \textbf{normas del examen}.
    Si se han recuperado normas de una evaluación
    anterior de la misma asignatura, compruebe que
    siguen siendo adecuadas.

    \item Pulse \textbf{Continuar}.
\end{enumerate}

\noindent\textbf{Validaciones que aplica el sistema}

\begin{itemize}
    \item La fecha y la hora son obligatorias.
    \item No se admite programar una nueva evaluación
    en una fecha anterior a la actual.
    \item La hora de inicio debe encontrarse
    entre las 06:00 y las 22:00.
    \item La duración debe estar entre
    15 y 480 minutos.
    \item El examen no debe terminar después
    de medianoche.
    \item Deben respetarse las fechas del período
    y la ventana académica correspondiente.
\end{itemize}

Si algún dato incumple estas restricciones,
el sistema informa el motivo y mantiene al usuario
en la etapa para que lo corrija.

\CapturasRegistro
{1}
{Etapa 1, información general y normas}
{Los campos se organizan verticalmente.
Desplácese para revisar las normas y localizar
el botón Continuar. No avance sin comprobar
fecha, hora y duración.}

% ============================================================
% ETAPA 2
% ============================================================

\subsubsection{Etapa 2. Selección de grupos}

Un examen debe vincularse a los grupos académicos
que participarán. Esta selección determina parte
del alcance de estudiantes y las opciones posteriores
de asignación de ambientes.

\noindent\textbf{Procedimiento}

\begin{enumerate}
    \item Compruebe que aparece la asignatura
    seleccionada en la primera etapa.

    \item Localice los grupos académicos disponibles.

    \item Seleccione los grupos que rendirán
    la evaluación.

    \item Si hay numerosos grupos, utilice
    el buscador disponible para encontrarlos.

    \item Revise la selección antes de continuar.
\end{enumerate}

\noindent\textbf{Comprobación}

El formulario debe conservar al menos un grupo
seleccionado. El código del grupo y la información
docente permiten comprobar que la selección
corresponde a los participantes previstos.

No seleccione grupos ajenos al examen únicamente
para avanzar en el formulario.

\CapturasRegistro
{2}
{Etapa 2, selección de grupos académicos}
{Revise las casillas de selección.
En pantallas estrechas puede ser necesario
desplazarse para consultar el listado completo
y las acciones de navegación.}

% ============================================================
% ETAPA 3
% ============================================================

\subsubsection{Etapa 3. Asignación de aulas}

En esta etapa se eligen los ambientes en los
que se realizará la evaluación. Las opciones
disponibles se consultan tomando en cuenta
los grupos seleccionados.

\noindent\textbf{Procedimiento}

\begin{enumerate}
    \item Abra la etapa \textbf{Aulas}.

    \item Examine los ambientes que presenta
    el sistema.

    \item Identifique cada aula por su nombre
    y la ubicación informada.

    \item Seleccione los ambientes necesarios
    para organizar el examen.

    \item Compruebe que la selección coincide
    con la planificación de ambientes.

    \item Continúe a la asignación de personal.
\end{enumerate}

\noindent\textbf{Comprobación}

Debe existir al menos un ambiente seleccionado
para continuar con un examen que requiere control
de ingreso por aula.

\noindent\textbf{Precaución}

La selección visual no sustituye la revisión de
la planificación académica ni la comprobación
de restricciones que informe el servidor.

\CapturasRegistro
{3}
{Etapa 3, selección de ambientes}
{La vista reorganiza los controles de aulas.
Compruebe los nombres de los ambientes seleccionados;
si existe un panel lateral, desplácese para consultarlo.}

% ============================================================
% ETAPA 4
% ============================================================

\subsubsection{Etapa 4. Personal responsable}

Cada aula seleccionada debe relacionarse con
la información del personal que atenderá
el examen. El formulario contempla una persona
encargada y personal de apoyo cuando corresponda.

\noindent\textbf{Procedimiento}

\begin{enumerate}
    \item Acceda a \textbf{Quién atiende}.

    \item Identifique las aulas que fueron
    seleccionadas previamente.

    \item Revise la persona encargada de
    cada ambiente.

    \item Compruebe la asignación del personal
    de apoyo correspondiente.

    \item Verifique que no existan asignaciones
    incorrectas antes de registrar.

    \item Pulse \textbf{Registrar examen}
    si se trata de uno nuevo, o
    \textbf{Guardar examen} si está editando
    un registro propio existente.
\end{enumerate}

\noindent\textbf{Importante}

Al guardar, la aplicación envía los datos al
servidor. Si este rechaza algún campo, el
formulario vuelve a la etapa donde debe
corregirse el problema.

Esta operación modifica información del sistema;
no debe ejecutarse únicamente para obtener
capturas de documentación.

\CapturasRegistro
{4}
{Etapa 4, encargados y personal de apoyo}
{Revise la correspondencia entre cada aula
y sus responsables. La acción de guardar
debe utilizarse solamente cuando la información
haya sido verificada.}

% ============================================================
% ETAPA 5
% ============================================================

\subsubsection{Etapa 5. Confirmación y resumen}

Cuando el servidor registra o actualiza
correctamente el examen, la aplicación muestra
la etapa \textbf{Listo}, con el resumen de
los datos guardados.

El usuario debe revisar:

\begin{enumerate}
    \item La asignatura y el tipo de evaluación.
    \item La fecha, hora y duración.
    \item Los grupos vinculados.
    \item Los ambientes asignados.
    \item El personal responsable de cada aula.
    \item Las normas configuradas para el examen.
\end{enumerate}

\noindent\textbf{Resultado esperado}

El sistema confirma el registro y permite
consultar la información guardada.

\noindent\textbf{Si aparece un error}

No considere completada la operación
hasta que el sistema confirme el guardado.
Regrese al campo señalado, corrija el dato
y vuelva a comprobar la respuesta.

\medskip
\noindent
{\color{Teal}\bfseries Continuación del proceso.}
Una vez creado el examen, la preparación
puede continuar con la habilitación de
estudiantes y los códigos QR, siempre que
la cuenta disponga de los permisos necesarios.

\begin{notebox}
\textbf{Cobertura visual de este procedimiento.}

Las figuras anteriores corresponden a
las cuatro etapas editables de un examen
existente, capturadas sin guardar cambios.
La confirmación de una operación nueva
debe verificarse funcionalmente antes de
añadir su captura como evidencia de registro.
\end{notebox}

% FIN DEL CAPÍTULO DE REGISTRO


\modulo{Mis grupos}{Docente}{/mis-grupos}{grupos}

La vista presenta grupos asociados al docente
y la lista de estudiantes correspondiente al grupo
seleccionado. Los registros deben consultarse
dentro del período pertinente.

\pasos
\begin{enumerate}
\item Abra \textbf{Mis grupos}.
\item Busque o seleccione la asignatura.
\item Seleccione el grupo deseado.
\item Consulte la lista de estudiantes y verifique
los datos académicos mostrados.
\item Si existe una opción de edición,
compruebe la autorización antes de utilizarla.
\end{enumerate}

\resultado{
El grupo activo coincide con el listado
que está siendo consultado.
}

\Pantallas
{docente/escritorio/grupos.png}
{docente/movil/grupos.png}
{2846}
{Mis grupos: selector de asignatura
y detalle del estudiantado.}

% ------------------------------------------------------------

\modulo{Mis reportes}
{Docente}{/mis-reportes}{misreportes}

A diferencia del reporte institucional, esta sección
presenta los reportes asociados a los exámenes del docente.
Puede incluir indicadores de habilitados, ausentes,
incidencias o intentos según el examen y el estado registrado.

\pasos
\begin{enumerate}
\item Abra \textbf{Mis reportes}.
\item Seleccione el examen correspondiente.
\item Revise los indicadores y categorías visibles.
\item Active los filtros de reporte pertinentes.
\item Descargue únicamente la información
permitida por su cuenta.
\end{enumerate}

\resultado{
La consulta y los archivos descargados corresponden
al examen seleccionado.
}

\Pantallas
{docente/escritorio/mis-reportes.png}
{docente/movil/mis-reportes.png}
{1856}
{Reportes del docente: selección de examen,
indicadores y descargas.}

% ------------------------------------------------------------

\modulo{Exámenes en curso}
{Docente}{/en-curso}{encurso}

La pantalla de seguimiento muestra indicadores del examen,
distribución por aulas y opciones relacionadas con ingreso
e incidencias. Los números reflejan el estado de la
evaluación en ese momento.

\pasos
\begin{enumerate}
\item Abra \textbf{En curso} desde el menú docente.
\item Identifique el examen activo.
\item Revise los indicadores de ingresos, faltantes,
aulas e incidencias.
\item Consulte el personal asignado a cada ambiente.
\item Abra el detalle de aula o incidencia solo
si requiere una actuación autorizada.
\end{enumerate}

\resultado{
El docente obtiene una vista consolidada del progreso
sin confundir conteos globales con un aula concreta.
}

\Pantallas
{docente/escritorio/en-curso.png}
{docente/movil/en-curso.png}
{1173}
{En curso: indicadores, responsables,
aulas y vista de ubicación.}

% ------------------------------------------------------------

\modulo{Control}{Docente}{/control}{control}

El punto de control muestra la información del examen
y del ambiente asignado cuando existe una sesión activa.
Las opciones de lectura o verificación dependen del estado
del examen y de los permisos.

\pasos
\begin{enumerate}
\item Abra \textbf{Punto de control}.
\item Confirme el examen y aula asignados.
\item Compruebe la disponibilidad del método de identificación.
\item Antes de confirmar un ingreso, verifique la habilitación
y el resultado presentado por el sistema.
\item Si no hay un examen activo,
no intente producir registros ficticios.
\end{enumerate}

\resultado{
El panel muestra controles acordes con la asignación
o informa que no existe una operación habilitada.
}

\Pantallas
{docente/escritorio/control.png}
{docente/movil/control.png}
{42}
{Punto de control docente: controles visibles
según el examen y su estado.}

% ============================================================
\section{Auxiliar: operación del punto de ingreso}
% ============================================================

\modulo{Punto de control}
{Auxiliar}{/ingreso}{ingreso}

El auxiliar trabaja exclusivamente sobre los exámenes
y ambientes asignados. En la evidencia recopilada puede
aparecer un estado vacío: no significa un error del sistema,
sino ausencia de una sesión abierta para esa cuenta.

\pasos
\begin{enumerate}
\item Acceda con la cuenta auxiliar.
\item Abra \textbf{Punto de control}.
\item Confirme si existe un examen abierto
en el aula asignada.
\item Cuando el sistema habilite el control,
siga el método de verificación y examine
la respuesta antes de autorizar un ingreso.
\item Si aparece un rechazo o incidencia,
remítalo al procedimiento institucional correspondiente.
\end{enumerate}

\resultado{
Si no hay examen asignado, la interfaz lo indica
y no permite registrar ingresos fuera de contexto.
}

\Pantallas
{auxiliar/escritorio/punto-control.png}
{auxiliar/movil/punto-control.png}
{0}
{Punto de control auxiliar: estado inicial condicionado
por las asignaciones.}

% ------------------------------------------------------------

\modulo{Mi aula}{Auxiliar}{/mi-aula}{miaula}

El módulo de lista del aula permite consultar
el estado de la asistencia. Su disponibilidad depende
de la ventana operativa del examen; la captura de referencia
muestra explícitamente un estado sin examen abierto.

\pasos
\begin{enumerate}
\item Abra \textbf{Mi aula}.
\item Revise el mensaje de estado que presenta el sistema.
\item Si hay examen activo, consulte la lista
y la información disponible para el aula.
\item Si no aparece ningún examen, verifique el período
y la asignación antes de solicitar soporte.
\end{enumerate}

\resultado{
La pantalla informa si existe una lista disponible
o si aún no se ha abierto el examen del aula.
}

\Pantallas
{auxiliar/escritorio/mi-aula.png}
{auxiliar/movil/mi-aula.png}
{0}
{Mi aula: la consulta de asistencia puede presentar
un estado vacío legítimo.}

% ============================================================
\section{Guía de resolución de incidencias habituales}
% ============================================================

{\small
\renewcommand{\arraystretch}{1.25}

\begin{tabularx}{\linewidth}{
 >{\raggedright\arraybackslash\bfseries}p{46mm}
 >{\raggedright\arraybackslash}X
}
\toprule
Situación & Comprobación y acción recomendada\\
\midrule

No aparece una opción &
Verifique el rol, los permisos efectivos y la asignación
académica; no suponga que todas las cuentas del rol poseen
el mismo acceso.\\

Examen no localizado &
Confirme el período, asignatura, filtros y fecha;
contraste con el docente responsable.\\

Punto de control vacío &
Revise si existe examen activo y aula asignada;
no genere datos de prueba en producción.\\

Error al exportar &
Revise filtros, sesión, conexión y permisos.
Informe el mensaje exacto al equipo responsable.\\

Registro o cifra inesperada &
Verifique el examen seleccionado, el alcance
de los filtros y la actualización de la consulta.\\

Respaldo fallido &
Documente el estado de ejecución y derive
el incidente a personal técnico autorizado.\\

\bottomrule
\end{tabularx}
}

% ============================================================
\section{Cobertura y límites de esta edición}
% ============================================================

Las \textbf{32 capturas} incluidas documentan
16 superficies verificadas, cada una en escritorio y móvil.

No se dispone todavía de imágenes del flujo estudiantil
ni de las etapas internas de creación de examen,
habilitación, emisión de QR y registro efectivo de ingresos;
por tanto, esas operaciones requieren capturas y validaciones
adicionales antes de declarar completo el manual operativo.

\begin{notebox}
\textbf{Control de calidad documental.}
Una captura de pantalla acredita la interfaz visible,
pero no sustituye una prueba de éxito de un procedimiento.
Al actualizar esta edición, cada paso crítico debe vincularse
con evidencia funcional y con el estado de la versión desplegada.
\end{notebox}

\end{document}
